\chapter{โครงข่ายไร้สาย LoRaWAN\index{โครงข่ายลอราแวน (LoRaWAN)}\index{LoRaWAN (AS923)} พลังงานแสงอาทิตย์และคลาวด์เทเลเมทรี}
\label{chap:ch09}

\section*{แผนบริหารการสอนประจำบทที่ 9}
\addcontentsline{toc}{section}{แผนบริหารการสอนประจำบทที่ 9}

\begin{planbox}{หัวข้อเนื้อหาประจำบท}
\begin{enumerate}[topsep=2pt, itemsep=1pt, partopsep=0pt, parsep=0pt, leftmargin=1.5em]
    \item ความท้าทายของการสื่อสารไร้สายข้ามหุบเขาและผืนน้ำไอเค็มทะเลในภาคตะวันออก
    \item สถาปัตยกรรมโครงข่ายระยะสั้น ESP-NOW และเกตเวย์สถานีฐานสองบอร์ด
    \item มาตรฐานโครงข่ายไร้สายระยะไกล LoRaWAN ย่านความถี่ AS923 และข้อกำหนดทางกฎหมาย
    \item เทคนิคการบีบอัดแพ็กเกตข้อมูลไบนารีระดับต่ำเพื่อลดระยะเวลาแพร่สัญญาณในอากาศ
    \item การออกแบบระบบพลังงานแสงอาทิตย์ร่วมกับแบตเตอรี่ LiFePO4\index{แบตเตอรี่ลิเธียมไอรอนฟอสเฟต (LiFePO4)}\index{LiFePO4 Battery} สำรองไฟ 3 วันสำหรับงานไอเค็ม
\end{enumerate}
\end{planbox}

\begin{planbox}{วัตถุประสงค์เชิงพฤติกรรม}
\begin{enumerate}[topsep=2pt, itemsep=1pt, partopsep=0pt, parsep=0pt, leftmargin=1.5em]
    \item อธิบายข้อแตกต่างระหว่างโหมด ESP-NOW, LoRa Point-to-Point, และ LoRaWAN ได้
    \item ออกแบบโครงสร้างข้อมูลไบนารีเพื่อบีบอัดค่าเซนเซอร์เกษตรให้มีขนาดไม่เกิน 12 ไบต์ได้
    \item คำนวณขนาดแผงโซลาร์เซลล์และความจุแบตเตอรี่ LiFePO4 ตามเกณฑ์สำรองไฟ 3 วันได้
    \item พัฒนาระบบเกตเวย์รับส่งข้อมูลผ่านโปรโตคอล MQTT TLS และแจ้งเตือนภัยผ่าน LINE ได้
\end{enumerate}
\end{planbox}

\newpage

\begin{planbox}{วิธีสอนและกิจกรรมการเรียนการสอน}
\begin{enumerate}[topsep=2pt, itemsep=1pt, partopsep=0pt, parsep=0pt, leftmargin=1.5em]
    \item บรรยายทฤษฎีการแพร่กระจายคลื่นวิทยุ แนวสายตา (Line-of-Sight) และโซนเฟรเนล (Fresnel Zone)
    \item สาธิตการเชื่อมต่อโมดูล LoRa ข้ามระยะทางจริง $2 - 5\text{ กิโลเมตร}$ ระหว่างแปลงเกษตร
    \item ปฏิบัติการคำนวณและประกอบชุดจ่ายไฟพลังงานแสงอาทิตย์พร้อมวงจรชาร์จ BMS
    \item การตั้งค่าเซิร์ฟเวอร์คลาวด์ ChirpStack และการสร้างแดชบอร์ดแสดงผลบน Grafana
\end{enumerate}
\end{planbox}

\begin{planbox}{สื่อการเรียนการสอน}
\begin{enumerate}[topsep=2pt, itemsep=1pt, partopsep=0pt, parsep=0pt, leftmargin=1.5em]
    \item เอกสารคำสอนบทที่ 9 สไลด์บรรยาย และคู่มือมาตรฐาน LoRaWAN AS923
    \item บอร์ด ESP32 พัฒนาพร้อมโมดูล SX1262 LoRa 923 MHz และเสาอากาศไฟเบอร์กลาส 5.8 dBi
    \item แผงโซลาร์เซลล์ชนิด Mono Half-Cut ขนาด 50W พร้อมแบตเตอรี่ LiFePO4 12V 20Ah
    \item ซอฟต์แวร์จำลองเครือข่ายและการรับส่งข้อความ MQTT ผ่าน Node-RED
\end{enumerate}
\end{planbox}

\begin{planbox}{การวัดและประเมินผล}
\begin{enumerate}[topsep=2pt, itemsep=1pt, partopsep=0pt, parsep=0pt, leftmargin=1.5em]
    \item การทดสอบความรู้เรื่องกฎหมายความถี่ กสทช. Duty Cycle และการคำนวณ Time-on-Air
    \item การประเมินความถูกต้องของโค้ดการเข้ารหัสและถอดรหัสข้อมูลไบนารี
    \item การทดสอบประสิทธิภาพการรับส่งข้อมูลจริงในพื้นที่สนามและความเสถียรของระบบคลาวด์
\end{enumerate}
\end{planbox}

\clearpage

\section{ความท้าทายของการสื่อสารไร้สายในภูมิประเทศภาคตะวันออก}

การติดตั้งระบบโทรมาตร (Telemetry) ในแปลงเกษตรกรรมและฟาร์มสัตว์น้ำของภาคตะวันออกต้องเผชิญกับอุปสรรคทางกายภาพที่แตกต่างจากพื้นที่ทั่วไป
\begin{itemize}[leftmargin=1.5em]
    \item \textbf{สวนผลไม้บนพื้นที่ลาดชันและหุบเขา (เช่น เขาคิชฌกูฏ โป่งน้ำร้อน สอยดาว):} มีแนวเขาสูงและต้นไม้ทึบบดบังสัญญาณโทรศัพท์เคลื่อนที่ สัญญาณ Wi-Fi ธรรมดาไม่สามารถทะลุผ่านมวลใบไม้หนาแน่นได้เกิน $50 - 100\text{ เมตร}$
    \item \textbf{ฟาร์มเพาะเลี้ยงสัตว์น้ำชายฝั่ง (เช่น แหลมสิงห์ ขลุง แหลมงอบ คลองใหญ่):} ผืนน้ำขนาดกว้างใหญ่ก่อให้เกิดการสะท้อนของคลื่นวิทยุหลายเส้นทาง (Multipath Fading) และมีละอองไอเค็มทะเลเข้มข้น ซึ่งกัดกร่อนขั้วต่อเสาอากาศและแผงวงจรอิเล็กทรอนิกส์อย่างรวดเร็ว
\end{itemize}

\section{สถาปัตยกรรมโครงข่ายไร้สายผสมผสาน}

เพื่อความคุ้มค่าและเสถียรภาพสูงสุด ระบบได้จัดแบ่งการสื่อสารออกเป็น 2 ลำดับชั้น ได้แก่ โครงข่ายระยะใกล้ในแปลงด้วย \textbf{ESP-NOW} และโครงข่ายข้ามแปลงระยะไกลด้วย \textbf{LoRaWAN}

\begin{table}[H]
\centering
\caption{การเปรียบเทียบเทคโนโลยีการสื่อสารไร้สายสำหรับงานเกษตรกรรมและสัตว์น้ำ}
\label{tab:wireless_comparison}
\begin{tabularx}{\textwidth}{p{3.2cm}YYYY}
    \toprule
    \textbf{คุณสมบัติ} & \textbf{ESP-NOW} & \textbf{LoRa (Raw P2P)} & \textbf{LoRaWAN} & \textbf{Cellular 4G/NB-IoT} \\
    \midrule
    ระยะรับส่งสัญญาณ & 100 - 300 เมตร & 2 - 10 กิโลเมตร & 5 - 15 กิโลเมตร & ทั่วประเทศที่มีสัญญาณ \\
    ขนาด Payload & สูงสุด 250 ไบต์ & ไม่เกิน 250 ไบต์ & 51 - 222 ไบต์ & ไม่จำกัด \\
    การพึ่งพาอินเทอร์เน็ต & ไม่ต้องต่อเน็ต & ไม่ต้องต่อเน็ต & ต้องต่อที่ Gateway & ต้องมีซิมการ์ดทุกโหนด \\
    อัตรากินพลังงาน & ปานกลาง (Wi-Fi PHY) & ต่ำมาก & ต่ำมาก (แบตเตอรี่อยู่ได้ปี) & สูง (ต้องชาร์จไฟบ่อย) \\
    ค่าใช้จ่ายรายเดือน & ฟรี ไม่มีค่าบริการ & ฟรี ไม่มีค่าบริการ & ฟรี (เครือข่ายส่วนตัว) & มีค่าบริการซิมการ์ด \\
    \bottomrule
\end{tabularx}
\end{table}

\section{มาตรฐาน LoRaWAN AS923 และการบีบอัดไบนารี}

ตามประกาศของสำนักงาน กสทช. ประเทศไทยกำหนดให้ใช้งานคลื่นความถี่ LoRa ในย่าน \textbf{920 ถึง 925 MHz (AS923)} โดยมีข้อกำหนดทางกฎหมายที่เคร่งครัดว่า อุปกรณ์จะต้องมีสัดส่วนเวลาส่งคลื่นวิทยุออกสู่อากาศ \textbf{Duty Cycle ไม่เกินร้อยละ 1} หรือหมายความว่าใน 1 ชั่วโมง อุปกรณ์สามารถแพร่สัญญาณรวมกันได้ไม่เกิน \textbf{36 วินาที}

เพื่อปฏิบัติตามกฎหมายและประหยัดพลังงานแบตเตอรี่ ข้อมูลการตรวจวัดทั้งหมด 7 ค่า (ดินหรือน้ำ) จะต้องถูกบีบอัดเป็นรหัสไบนารีขนาดกะทัดรัดเพียง \textbf{11 ไบต์} ดังโครงสร้างในโค้ดต่อไปนี้

\begin{pythonbox}[การบีบอัดข้อมูลเซนเซอร์ 7 ค่าลงในไบนารี 11 ไบต์ด้วย Python]
import struct

def pack_aqua_lora_payload(do_mg_l: float, temp_c: float, salinity_ppt: float, 
                           ph: float, battery_v: float) -> bytes:
    # 1. DO (0 - 20.00 mg/L) คูณ 100 เก็บเป็น unsigned short (2 bytes)
    do_raw = int(do_mg_l * 100)
    # 2. อุณหภูมิ (-10.0 ถึง 60.0 C) คูณ 10 เก็บเป็น signed short (2 bytes)
    temp_raw = int(temp_c * 10)
    # 3. ความเค็ม (0 - 50.0 ppt) คูณ 10 เก็บเป็น unsigned short (2 bytes)
    sal_raw = int(salinity_ppt * 10)
    # 4. pH (0 - 14.00) คูณ 100 เก็บเป็น unsigned short (2 bytes)
    ph_raw = int(ph * 100)
    # 5. แรงดันแบตเตอรี่ (0 - 15.00 V) คูณ 100 เก็บเป็น unsigned short (2 bytes)
    bat_raw = int(battery_v * 100)
    # 6. บิตสถานะเตือนภัย (Flags) เก็บเป็น unsigned char (1 byte)
    status_flags = 0x01 if do_mg_l < 4.0 else 0x00
    
    # รวมข้อมูลไบนารี Big-Endian ขนาดรวม 2+2+2+2+2+1 = 11 Bytes พอดี
    payload = struct.pack(">HhHHHB", do_raw, temp_raw, sal_raw, ph_raw, bat_raw, status_flags)
    return payload
\end{pythonbox}

\section{การออกแบบระบบพลังงานแสงอาทิตย์ทนไอเค็มทะเล}

โหนดตรวจวัดกลางแจ้งและสถานีตรวจวัดคุณภาพน้ำกลางบ่อกุ้งต้องพึ่งพาพลังงานแสงอาทิตย์แบบพึ่งพาตนเอง 100\% การคำนวณขนาดแผงโซลาร์เซลล์และแบตเตอรี่อิงเกณฑ์ \textbf{การสำรองไฟในสภาวะฝนตกต่อเนื่อง 3 วัน (3 Days of Autonomy)}

\begin{equation}
E_{\text{daily}} = V_{\text{system}} \times I_{\text{average}} \times 24 \quad (\text{Wh/day})
\label{eq:daily_energy}
\end{equation}

\begin{equation}
C_{\text{battery}} = \frac{E_{\text{daily}} \times N_{\text{autonomy}}}{V_{\text{system}} \times \text{DOD} \times \eta_{\text{bms}}} \quad (\text{Ah})
\label{eq:battery_capacity}
\end{equation}

โดยที่ $N_{\text{autonomy}} = 3\text{ วัน}$, $\text{DOD} = 0.80$ (ความลึกของการคายประจุปลอดภัยของแบตเตอรี่ลิเธียมไอรอนฟอสเฟต LiFePO4), และ $\eta_{\text{bms}} = 0.95$

\begin{theorybox}[ข้อกำหนดการติดตั้งอุปกรณ์ในพื้นที่ไอเค็มชายฝั่ง]
อุปกรณ์ทุกชิ้นที่ติดตั้งในฟาร์มเพาะเลี้ยงสัตว์น้ำชายฝั่งต้องมีมาตรฐานการป้องกัน:
\begin{itemize}[leftmargin=1.5em]
    \item ตัวกล่องควบคุมต้องเป็นวัสดุโพลีคาร์บอเนตกันรังสี UV มาตรฐานป้องกันน้ำและฝุ่นระดับ \textbf{IP67}
    \item น็อตและขายึดโลหะทั้งหมดต้องเป็นสแตนเลสเกรดมารีน \textbf{SUS316L} ที่ทนต่อการกัดกร่อนของคลอไรด์
    \item แผงโซลาร์เซลล์ต้องผ่านการรับรองมาตรฐานต้านทานไอเค็มทะเล \textbf{IEC 61701 Salt Mist Corrosion Severity 6}
\end{itemize}
\end{theorybox}

\section*{คำถามทบทวนท้ายบทที่ 9}
\addcontentsline{toc}{section}{คำถามทบทวนท้ายบทที่ 9}

\begin{enumerate}
    \item จงอธิบายข้อกำหนด Duty Cycle ร้อยละ 1 ตามประกาศ กสทช. ย่านความถี่ AS923 และวิเคราะห์ว่าส่งผลต่อการออกแบบความถี่ในการส่งข้อมูลเซนเซอร์อย่างไร
    \item เหตุใดการบีบอัดข้อมูลด้วย Struct Packing เป็นรหัสไบนารี 11 ไบต์ จึงช่วยประหยัดแบตเตอรี่ของโหนด LoRaWAN ได้มากกว่าการส่งข้อมูลในฟอร์แมต JSON
    \item จงคำนวณหาความจุแบตเตอรี่ LiFePO4 12V ที่ต้องการ หากอุปกรณ์ตรวจวัดกินพลังงานเฉลี่ยวันละ 18 Wh ภายใต้เงื่อนไขสำรองไฟ 3 วัน และ DOD ร้อยละ 80
    \item เปรียบเทียบข้อดีและข้อจำกัดของการใช้โครงข่าย ESP-NOW เชื่อมต่อภายในสวน เทียบกับการใช้ LoRaWAN ทุกโหนด
\end{enumerate}
\clearpage
